\section{Companion lab: variable-length records in a ring}
\label{revision-lab:ring-buffer}

Life reuses an allocation after its old contents are no longer needed. A queue
reuses storage too, but it may contain several records with different lifetimes.
The oldest record must remain available until it is successfully removed, while
new records occupy the free space after the newest one. A circular buffer, or
\emph{ring buffer}, lets that occupied region cross the physical end of an array.
The array stays in place; its indices wrap back to zero.

This companion is original project \texttt{014-ring-buffer-logger}. Use the revised
\texttt{ring-records.asm} starter in \texttt{revision/labs/files} and both checks
under \texttt{packages/014-ring-buffer-logger/checks}. Implement the
callable \texttt{push\_record} and \texttt{pop\_record} functions. The supplied
caller, allocations and \texttt{operation\_ready} marker provide a common way to
observe and check your implementation. Replace the failing function bodies; choose your own register allocation and algorithm.
The archived project starter remains available under its original package path.

You will need the preceding chapters' byte addressing, function interfaces and
preserved state. The new work is to separate physical positions from logical
record order, account for prefix bytes, and arrange failure checks before changing
caller data. The complete queue functions are your implementation task; the short
cursor experiment below isolates one operation you can verify first.

\subsection*{Represent a queue with indices and an occupied count}

A record consists of one unsigned length byte followed by $L$ payload bytes, where
$0\leq L\leq31$. The payload is binary data: zero is an ordinary byte, not a string
terminator. An empty payload is a valid record and still occupies its one-byte
prefix. The physical ring capacity $C$ is from one through 32 bytes and need not
be a power of two.

Each queue has an eight-byte-aligned, 24-byte descriptor. Offsets are byte
displacements from the descriptor pointer passed in \texttt{\$a0}.

\begin{center}
\begin{tabular}{rll}
\toprule
Offset & Representation & Meaning\\
\midrule
0 & word & head: first byte of the oldest record\\
4 & word & tail: position for the next appended prefix\\
8 & word & used: occupied bytes, including prefixes\\
12 & word & capacity: usable ring bytes\\
16 & doubleword & full address of the ring allocation\\
\bottomrule
\end{tabular}
\end{center}

A valid descriptor satisfies $0\leq head,tail<C$, $0\leq used\leq C$, and
\[
 tail=(head+used)\bmod C.
\]
The occupied circular bytes contain complete valid records. You may assume this
input invariant; this lab does not ask you to repair corrupted descriptors.
The supplied descriptor, ring, source and output are separate objects. Data pointers
may be byte unaligned. Load the ring pointer with \texttt{ld} from offset 16, then
use full-width address arithmetic when adding an index. Loading only its low word
would change the interface even if a particular Default-layout address happened
to fit in that word.

Head equal to tail is insufficient to identify an empty queue. An empty queue has
$used=0$; a full queue has $used=C$. Both have equal head and tail. The occupied
count also makes every ring byte usable, rather than reserving a slot to distinguish
those states. Bytes outside the occupied region can retain stale data. Their
values do not make them live records.

\subsection*{Follow a record across the physical boundary}

Start with capacity seven and an empty queue whose head and tail are both five.
Push the three-byte payload \texttt{ABC}, whose byte values are 65, 66 and 67.
Its prefix is three, so it occupies four positions: 5, 6, 0 and 1. Tail becomes
two and used becomes four. The next two-byte payload \texttt{xy} occupies positions
2, 3 and 4, including its prefix. Tail returns to five and used becomes seven.
Figure~\ref{fig:revision-ring-full} shows the resulting full queue.

\begin{figure}[htbp]
  \centering
  \begin{tikzpicture}[x=10mm,y=10mm,alt={A full seven-byte ring has head and tail both five and used seven. Physical indices zero through six contain decimal bytes 66, 67, 2, 120, 121, 3, 65. The oldest record begins with length three at index five, then payload A at six, B at zero and C at one. The second record has length two at index two and payload x and y at three and four. Shading marks the first record; a wrap arrow connects the physical ends.}]
    \foreach \i in {0,1,5,6} {\fill[black!18] (\i,0) rectangle ++(1,1);}
    \foreach \i in {0,...,7} {\draw (\i,0) -- (\i,1);}
    \draw (0,0) -- (7,0) (0,1) -- (7,1);
    \foreach \i/\value/\kind in {0/66/B,1/67/C,2/2/length,3/120/x,4/121/y,5/3/length,6/65/A} {
        \node[font=\small,anchor=south] at (\i+0.5,1.05) {\i};
        \node[font=\small] at (\i+0.5,0.65) {\value};
        \node[font=\footnotesize] at (\i+0.5,0.25) {\kind};
      }
    \draw[-{Stealth},rounded corners=2pt] (7,0.5) -- (7.4,0.5) -- (7.4,1.7) -- (-0.4,1.7) -- (-0.4,0.5) -- (0,0.5);
    \draw[-{Stealth}] (5.5,-0.7) -- (5.5,0);
    \node[anchor=north,font=\small] at (5.5,-0.75) {$head=tail=5$};
    \node[anchor=north,font=\small] at (2,-0.2) {$used=7$ (full)};
  \end{tikzpicture}
  \caption{Logical record order crosses the physical boundary. Each length prefix occupies a byte of capacity, and equal head and tail indices here describe a full queue.}
  \label{fig:revision-ring-full}
\end{figure}


Read the oldest record in logical order, beginning at head, rather than by scanning
physical index zero upward. A pop with destination capacity two cannot hold
\texttt{ABC}. It returns a failure and the required payload length three. Nothing
is consumed, so a retry with capacity three must return the same record. After
that successful pop, head becomes two and used becomes three; tail remains five.
The second record is still queued.

\begin{center}
\begin{tabular}{lrrrl}
\toprule
Completed operation & Head & Tail & Used & $(v_0,v_1)$\\
\midrule
Initial empty queue & 5 & 5 & 0 & ---\\
Push \texttt{ABC} & 5 & 2 & 4 & $(0,4)$\\
Push \texttt{xy} & 5 & 5 & 7 & $(0,3)$\\
Pop, capacity 2 & 5 & 5 & 7 & $(-2,3)$\\
Pop, capacity 3 & 2 & 5 & 3 & $(3,4)$\\
Pop, capacity 31 & 5 & 5 & 0 & $(2,3)$\\
\bottomrule
\end{tabular}
\end{center}

Figure~\ref{fig:revision-ring-memory} shows those full and empty descriptors in the
native Memory tab for this seven-byte example. Each group of four hexadecimal
bytes is one little-endian word; the boxed third group holds the occupied count.

\begin{figure}[htbp]
\centering
\begin{tikzpicture}[alt={Two actual native Memory tab observations of the seven-byte example. The first descriptor row contains head five, tail five, used seven and capacity seven: it is full. The second contains head five, tail five, used zero and capacity seven: it is empty. Native publication boxes mark the third four-byte group, which changes from hexadecimal 07 00 00 00 to 00 00 00 00.}]
\node[anchor=north west,inner sep=0,font=\small\bfseries] (fulltitle) at (0,0) {1. Full queue};
\node[anchor=north west,inner sep=0] (full) at ([yshift=-2mm]fulltitle.south west) {\MARSRevisionCrop{revision/captures/ring/ring-full.png}{0.94\linewidth}{0.034375000}{0.337777778}{0.584375000}{0.608888889}};
\node[anchor=north west,inner sep=0,font=\small\bfseries] (emptytitle) at ([yshift=-7mm]full.south west) {2. Empty queue};
\node[anchor=north west,inner sep=0] (empty) at ([yshift=-2mm]emptytitle.south west) {\MARSRevisionCrop{revision/captures/ring/ring-empty.png}{0.94\linewidth}{0.034375000}{0.337777778}{0.584375000}{0.608888889}};
\foreach \panel in {full,empty} {
  \coordinate (\panel-left) at ($(\panel.south west)!0.586065574!(\panel.south east)$);
  \coordinate (\panel-right) at ($(\panel.south west)!0.774590164!(\panel.south east)$);
  \coordinate (\panel-bottom) at ($(\panel.south west)!0.020833333!(\panel.north west)$);
  \coordinate (\panel-top) at ($(\panel.south west)!0.458333333!(\panel.north west)$);
  \draw[black,line width=0.7pt] (\panel-left |- \panel-bottom) rectangle (\panel-right |- \panel-top);
}
\end{tikzpicture}
\caption{Actual light-theme Debug observations at completed commands in the seven-byte example. Head and tail remain five in both rows; the boxed used word distinguishes full from empty. The four words, from left to right, are head, tail, used and capacity.}
\label{fig:revision-ring-memory}
\end{figure}


The final empty state has the same indices as the full state. Popping does not
clear ring bytes: the occupied count and head release their ownership. The old
bytes remain available for a later push to overwrite. This also means that a
successful pop writes its destination and descriptor, but never the ring itself.

\subsection*{Establish the cursor operation separately}

For a one-byte advance, increment the index and replace it with zero when it
reaches capacity. Since the old index is at most $C-1$, the increment is at most
$C$, so an equality check is sufficient. The complete \texttt{ring-cursor.asm}
experiment in \texttt{revision/labs/files} calls this helper with index six and
capacity seven, retains the result in \texttt{cursor\_result}, and exits silently.
The helper receives the index in \texttt{\$a0}, capacity in \texttt{\$a1}, and
returns the next index in \texttt{\$v0}.

\noindent\begin{minipage}{\linewidth}
\MARSCodeLabel{Advance one byte without assuming a power-of-two capacity}
\begin{minted}{mars-idea}
advance_cursor:
    addiu $v0, $a0, 1
    bnec $v0, $a1, cursor_return
    nop
    move $v0, $zero
cursor_return:
    jrc $ra
\end{minted}
\end{minipage}

Run its public \texttt{revision/labs/checks/ring-cursor.kts} check to cover every
valid index for capacities one through 32. Predict capacity one before running:
its only valid index advances back to zero. Then compare a mask-based expression
$(index+1)\mathbin{\&}(C-1)$. For $C=7$ and index six, it produces six instead of
zero. Keeping selected low bits implements a power-of-two modulus, so that shortcut
does not satisfy the arbitrary-capacity contract.

The equality shortcut above belongs to a one-byte increment. A whole-record jump
can pass capacity without equaling it. In the example, $5+4=9$, whose wrapped value
is two. You can advance the cursor once per copied byte, use a remainder operation,
or subtract $C$ when a bounded whole-record sum is at least $C$. State the bound
that makes your chosen method valid. A fitting record occupies at most $C$ bytes,
so its old index plus encoded size is less than $2C$.

If you call the cursor helper from a queue function, that function becomes a
non-leaf caller and must preserve its return link and live values across the call.
You may instead inline the same bounded decision. The grading contract accepts
either organization.

\subsection*{Validate the request before copying}

Both queue calls receive the descriptor in \texttt{\$a0}. For push,
\texttt{\$a1} points to $L$ source payload bytes and \texttt{\$a2} is $L$.
For pop, \texttt{\$a1} points to output storage and \texttt{\$a2} is that storage's
capacity, also from zero through 31. The two result registers have these meanings:

\begin{center}
\begin{tabular}{lll}
\toprule
Situation & $v_0$ & $v_1$\\
\midrule
Push: $L\geq C$ & $-2$ & 0\\
Push: insufficient free bytes & $-1$ & 0\\
Push: success & 0 & $L+1$\\
Pop: empty queue & $-1$ & 0\\
Pop: destination too small & $-2$ & required $L$\\
Pop: success & $L$ & $L+1$\\
\bottomrule
\end{tabular}
\end{center}

For push, first reject a record that cannot fit even in an empty ring: $L\geq C$.
Then compare its encoded size $L+1$ with the free count $C-used$. This order gives
the specified error priority when the ring is already full and the record is also
too large. A representable record with insufficient current space returns $-1$;
a record whose own size is impossible returns $-2$.

For pop, test $used=0$ before interpreting a prefix. An empty ring can contain stale
bytes that have no record meaning. When a record exists, read its length and compare
it with the destination capacity. If it does not fit, return $(-2,L)$ without
advancing head or writing any output. The caller can then retry with enough storage.

A failure permits no stores to caller data, including a temporary write that is
later restored. Arrange the complete failure decision before such stores. Registers
can hold candidate indices and counts while you check the request. A private callee
stack frame is separate from caller-owned descriptor, ring, source and output data;
restore the agreed saved registers and stack pointer on every return path.

For capacity one, a zero-length push succeeds: its prefix consumes the one available
byte. A pop into capacity zero then succeeds with $(0,1)$ and performs no output
store. Distinguish that zero-length success from an empty-queue failure. The second
return value accounts for the consumed prefix even when there is no payload.

\subsection*{Implement the two successful transitions}

For a successful push, begin at the old tail and store the length prefix. Advance
the working cursor once, then copy exactly $L$ source bytes, advancing and wrapping
the ring cursor after each byte. Source payload bytes may contain zero,
\texttt{0x80} or \texttt{0xff}; copy by length, using byte loads and stores.
A zero-length payload skips the copy loop and performs no source load. On completion,
retain the original head, publish the final tail, and increase used by $L+1$.
Capacity and the full ring pointer are read-only descriptor fields.

For a successful pop, retain the length from the validated prefix. Advance a working
cursor past it and copy exactly $L$ ring bytes to the destination. Do not copy the
prefix or append a terminator. Afterward publish the resulting head, subtract $L+1$
from used, and preserve tail. Ring bytes remain unchanged throughout the call. A
record can cross the boundary between any pair of payload bytes, including the
pair immediately after its prefix.

Keep a loop invariant beside each implementation. After $k$ payload bytes,
$0\leq k\leq L$, the working ring cursor has advanced $k+1$ positions from its
original head or tail; the extra one accounts for the prefix. The source or output
pointer has advanced only $k$ bytes. This distinction locates many off-by-one
faults without relying on a plausible final string.

You may choose the successful descriptor-store order and copy direction as long
as the final state and access bounds satisfy the contract. Keeping working state
in registers until the copy is complete makes the transition easier to inspect.
The supplied caller publishes results only after the function returns. At
\texttt{operation\_ready}, the descriptor, \texttt{last\_status},
\texttt{last\_aux} and \texttt{operation\_index} describe that completed command.
An intermediate instruction inside a call need not yet have that completed-state
interpretation.

The revised caller aligns the stack and uses the same compact call and branch
forms as the main book. Its unfinished functions stop with \texttt{break} until
you supply their bodies. The archived project starter also remains compatible
with these checks; it uses retained delayed transfers with explicit \texttt{nop}
slots. On this MIPS64 target, a retained \texttt{jal} executes its following slot
before entering the function and saves the call address plus eight. Its
\texttt{jr} slot executes before returning. Keep that distinction when reading
the archived source; the revised caller and cursor experiment use compact calls.

\subsection*{Test state, access bounds and recovery}

Begin with one record that fits without wrapping. Check its prefix, payload, updated
tail and occupied count. Pop it and check the exact output byte count, the released
space and the unchanged ring. Next reproduce the seven-byte example above, including
the failed short-destination pop and its retry. Write down all physical positions
before running it.

The supplied 14-command caller uses capacity eight and a different starting state.
Break at \texttt{operation\_ready} and use \texttt{operation\_index} to identify the
completed command; the command index starts at zero. Index three shows a wrapped
queue. Index eight is full. Index ten is a short-destination failure, and index
eleven retries with enough capacity. Compare the full ring and descriptor across
the failure, then locate the bytes released by the successful retry. Output bytes
after the requested payload must retain their previous values.

Run both public checks on your completed copy, using JDK 25:
\begin{verbatim}
java -jar /path/to/cli.jar grade \
  packages/014-ring-buffer-logger/checks/kernels.kts \
  /path/to/your/main.asm
java -jar /path/to/cli.jar grade \
  packages/014-ring-buffer-logger/checks/program.kts \
  /path/to/your/main.asm
\end{verbatim}
The kernel check supplies different valid descriptors, capacities, offsets, saved
register values and return continuations. The caller check observes completed
commands and result publication. They check byte footprints and protected stores,
not just a decoded payload string. An unchanged final descriptor does not excuse
an earlier forbidden store on a failed call.

For diagnosis, replace your wrap logic with a mask in a copy and use capacity seven.
Then try accounting for only $L$ occupied bytes, omitting the prefix, and use a
zero-length record. Finally, try advancing head before the pop capacity check and
repeat the failed-pop/retry trace. In each case, identify the first incorrect
state transition or store before repairing the code. Keep those counterexamples
with a passing report and an explanation of why empty and full can share indices.
